\scanpage{134}
% archive foliation: 133
i.e.
\figure{\begin{tikzpicture}\draw (0,0)--(0,.32);\end{tikzpicture}}
\diagram{\begin{tikzcd}[ampersand replacement=\&]
L_\omega \arrow[r, no head, "\sim"] \& V\\
z \arrow[r, mapsto] \& \operatorname{Re}z
\end{tikzcd}}
est un isomorphisme d'espaces vectoriels complexes, quand on munit $V$ de la structure complexe $c_\omega$. ]

\margin{Mais il y a lieu d'associer à cette structure complexe sur $V$ \unsure{non pas} $L_\omega$, mais $L_{\omega'}$, cf. plus bas\ldots}

Dans le cas qui nous occupe, $\operatorname{rg}_{\mathbb R}V=2$ les \deleted{droites} sous-espaces vectoriels envisagés de $V_{\mathbb C}$ sont de \deleted{$\unsure{\operatorname{rg}}$} $\operatorname{rg}1$, \inserted{$\operatorname{corg}1$,} et peuvent s'interpréter comme les points de
\begin{equation}
\Sigma=\Sigma_V=P(\mathbb C),\quad\text{où }P=\mathbb P(V)
\tag{16}
\end{equation}
est le \deleted{\unsure{esp.}} schéma droite projective réelle associé à $V$. Pour que
\diagram{\begin{tikzcd}[ampersand replacement=\&]
L \arrow[r] \& V\\
z \arrow[r, mapsto] \& \operatorname{Re}z
\end{tikzcd}}
soit un isomorphisme, i.e. pour que $L$ corresponde à une structure complexe \deleted{de} sur $V$, il faut que \deleted{le point $z$ correspondant de $\Sigma$ ait \unsure{une pensée|une partie}} \inserted{$L\cap V=\{0\}$ i.e. $L$} n'est pas complexifié d'une droite $L_0\subset V_{\unsure{\mathbb R}}$; i.e. que le point $z\in\Sigma$ qui correspond à $L$ \inserted{ne} soit pas dans $\deleted{\Sigma(\mathbb R)}$
\[
\Sigma_{\mathbb R}=P(\mathbb R)
\]
ou encore qu'on ait
\begin{equation}
z\in\Sigma\setminus\Sigma_{\mathbb R}.
\tag{17}
\end{equation}
\deleted{\unsure{En résumé}} Ainsi, les structures complexes
